% Generated draft: requires mathematical and source review.
\documentclass[11pt]{article}
\usepackage[margin=25mm]{geometry}
\usepackage{fontspec}
\setmainfont{Latin Modern Roman}
\usepackage{amsmath,amssymb,mathtools,unicode-math}
\setmathfont{Latin Modern Math}
\usepackage{enumitem,needspace,xurl,hyperref}
\setlength{\parindent}{0pt}
\setlength{\parskip}{5pt}
\setlength{\emergencystretch}{3em}
\newcommand{\N}{\mathbb{N}}
\newcommand{\Z}{\mathbb{Z}}
\newcommand{\Q}{\mathbb{Q}}
\newcommand{\R}{\mathbb{R}}
\newcommand{\C}{\mathbb{C}}
\newcommand{\F}{\mathbb{F}}
\newcommand{\A}{\mathbb{A}}
\newcommand{\PP}{\mathbb{P}}
\newcommand{\E}{\mathbb{E}}
\newcommand{\eps}{\varepsilon}
\newcommand{\dd}{\,\mathrm{d}}
\begin{document}
\section*{UnsolvedMath ID: 30002271}
\textbf{Source problem code:} catalogue entry
\subsection{Definitions and mathematical statement}
For context, under the usual convention that \(G=(V,E)\) is a nonempty finite simple undirected graph, the degree of \(u\in V\) is \(d_G(u)=|\{v\in V:\{u,v\}\in E\}|\). Let \(\mathcal I(G)\) be the family of independent vertex sets of \(G\). Its fractional chromatic number is
\[
\chi_f(G)=\min\left\{\sum_{I\in\mathcal I(G)}w_I:\ w_I\geq 0\text{ for every }I\in\mathcal I(G),\quad \sum_{\substack{I\in\mathcal I(G)\\v\in I}}w_I\geq 1\text{ for every }v\in V\right\}.
\]
The supplied record does not explicitly specify this graph domain; these definitions explain the standard finite-graph interpretation only.

The source question asks whether there exists a graph \(G\) such that
\[
\chi_f(G)>\max_{u\in V(G)}\frac{\displaystyle\sum_{v\in\widetilde N(u)}\gamma_\ell(v)}{d(u)+1}.
\]
Here \(d(u)\) is interpreted as the vertex degree. The supplied record does not define \(\widetilde N(u)\) or \(\gamma_\ell(v)\), including the meaning of the subscript \(\ell\). Consequently, the displayed existential question preserves the supplied expression but is not a fully specified mathematical proposition. The intended graph class and these two definitions must be recovered from the cited source before a precise formulation can be finalized. In particular, no identification of \(\widetilde N(u)\) with a closed neighborhood, or formula for \(\gamma_\ell(v)\), is asserted here.
\subsection{Short English statement}
The source asks whether some graph has fractional chromatic number strictly greater than the displayed maximum of local averages. The graph domain and the meanings of \(\widetilde N(u)\) and \(\gamma_\ell(v)\) are missing from the supplied record, so the question remains incompletely specified.
\subsection{Sources}
\begin{itemize}
\item Local UnsolvedMath registry, supplied complete record: permanent ID \(30002271\), source problem code \texttt{catalogue entry}, title ``Counterexamples to Averaged Local Fractional-Coloring Bounds,'' in the Oberwolfach Reports Open Problems collection. The record gives proposed year 2013 and names no proposer.
\item Supplied primary report citation: \emph{Graph Theory} (2013), Oberwolfach Reports, pp.~67--128, DOI 10.4171/owr/2013/02. \url{https://doi.org/10.4171/owr/2013/02}. The background describes this collection as drawn from workshop reports and their open-problem discussions. The registry labels the statement \texttt{corrected\_verified} and reports that its mathematical content was checked against the official report, with notation, context, and question form normalized without changing the source-level claim. This is supplied provenance, not an independent verification here.
\item Supplied literature source: Kelly--Postle, ``fractional coloring with local demands.'' \url{https://arxiv.org/abs/1811.11806}.
\item Supplied literature source: Edwards--King, ``local fractional relaxation.'' \url{https://doi.org/10.37236/2666}.
\item Supplied literature assessment, checked on \texttt{2026-08-21}: the registry labels the problem open and states that the stronger local fractional question remains conjectural in current primary work, with no general proof or counterexample found. This is supplied dated curation derived from the OWR collection, not independently verified here. The record explicitly calls for independent verification of time-sensitive status and cited claims.
\end{itemize}
\end{document}
